\documentclass[11pt,a4paper]{article}
\usepackage[margin=2.5cm]{geometry}
\usepackage[T1]{fontenc}
\usepackage{amsmath,amssymb}
\usepackage{amsthm}     % theorem environments and proof
\usepackage{tikz}       % drawing figures without external image files
\linespread{1.15}       % slightly more space between lines

% ---------------------------------------------------------------
% MODULE IV - Theorems, Figures, Cross References, Presentation
% Syllabus units 12-19 (Text 1 Tutorials IX-XIII; Text 2 Ch. 12).
% Compile twice with pdflatex (for \ref, \listoffigures).
% hyperref is loaded last among the packages, as recommended.
% ---------------------------------------------------------------

% --- theorem environments (amsthm) ---
\theoremstyle{plain}                         % bold heading, italic text
\newtheorem{theorem}{Theorem}[section]       % numbered within sections
\newtheorem{lemma}[theorem]{Lemma}           % shares the theorem counter
\theoremstyle{definition}                    % bold heading, upright text
\newtheorem{definition}[theorem]{Definition}
\theoremstyle{remark}                        % italic heading, upright text
\newtheorem{remark}[theorem]{Remark}

\usepackage[colorlinks=true,linkcolor=blue]{hyperref}   % load last

\title{Module IV: Theorems, Figures, Cross References and Presentation}
\author{Study Notes -- MAT5FS112}
\date{}

\begin{document}
\maketitle

\begin{abstract}
This note covers the AMS theorem package, boxes, floating figures,
cross references, footnotes and the layout conventions of a
presentable document. Each example compiles as it stands.
\end{abstract}

\tableofcontents
\listoffigures
\newpage

% ===============================================================
\section{Theorems with \texttt{amsthm}}
\label{sec:theorems}
% ===============================================================

\subsection{Defining and using a theorem}

The preamble defines environments with \verb|\newtheorem{name}{Title}|.
In the body the environment is used like any other:

\begin{theorem}[Pythagoras]\label{thm:pyth}
In a right-angled triangle with legs $a, b$ and hypotenuse $c$,
$a^2 + b^2 = c^2$.
\end{theorem}

\begin{lemma}
Every theorem in the same counter as Theorem~\ref{thm:pyth} is numbered
in sequence with it.
\end{lemma}

\begin{definition}
A triangle is \emph{right-angled} if one of its angles equals $90^\circ$.
\end{definition}

\begin{remark}
The theorem numbers include the section number because of the optional
argument \texttt{[section]} in \verb|\newtheorem|.
\end{remark}

\begin{proof}
Let the right angle be at $C$. Place a square of side $a+b$ around the
triangle; comparing areas gives the result. \qedhere
\end{proof}

\subsection{Choosing a theorem style}

\begin{verbatim}
\theoremstyle{plain}        % theorem, lemma, proposition, corollary
\theoremstyle{definition}   % definition, example, exercise
\theoremstyle{remark}       % remark, note, notation
\end{verbatim}

% ===============================================================
\section{Boxes}
% ===============================================================

\verb|\fbox| draws a frame around its contents. \verb|\framebox| takes
an optional width and alignment. \verb|\makebox| sets a fixed width
without a frame. \verb|\parbox| makes a fixed-width paragraph in a line.

\noindent
\fbox{Boxed text}\quad
\framebox[4cm][c]{Centred in 4\,cm}\quad
\makebox[3cm][r]{Right in 3\,cm}\\[1em]
\fbox{\parbox{0.7\textwidth}{A framed paragraph. The \texttt{parbox}
lets the text wrap inside a box of fixed width, for key results or
warnings.}}

% ===============================================================
\section{Floating figures}
% ===============================================================

A figure environment places the picture at a suitable location and
numbers it. Here the picture is drawn with TikZ so that the example
needs no image file. With an external picture, use
\verb|\includegraphics[width=0.6\textwidth]{file.png}| (\texttt{graphicx}).

\begin{figure}[htbp]
  \centering
  \begin{tikzpicture}[scale=1.2]
    \draw[->] (-0.3,0) -- (3.3,0) node[right] {$x$};
    \draw[->] (0,-0.3) -- (0,2.6) node[above] {$y$};
    \draw[thick, blue] (0,0) .. controls (1,2.2) and (2,2.2) .. (3,0.4);
    \draw[dashed] (0,0) -- (3,0.4);
    \node[below right] at (3,0.4) {end};
  \end{tikzpicture}
  \caption{A curve from the origin, with its chord drawn dashed.}
  \label{fig:curve}
\end{figure}

Figure~\ref{fig:curve} shows the curve. The caption sits below the
figure; \verb|\caption| before \verb|\label| is required so that the
reference shows the figure number.

% ===============================================================
\section{Cross references}
% ===============================================================

\subsection{Labels and the three reference commands}

A label is placed right after the item it names, inside its
environment or heading. Then:
\begin{itemize}
  \item \verb|\ref{label}| gives the number: Section~\ref{sec:theorems}.
  \item \verb|\pageref{label}| gives the page: page~\pageref{sec:theorems}.
  \item \verb|\eqref{label}| gives the equation number in brackets (for
    equations): \eqref{eq:mass}.
\end{itemize}

\begin{equation}\label{eq:mass}
  F = ma
\end{equation}

\subsection{Smart references with \texttt{hyperref}}

With \texttt{hyperref} loaded, \verb|\autoref{label}| writes the word
for you and makes the reference a clickable link: see
\autoref{fig:curve} and \autoref{thm:pyth}.

\subsection{Rules for labels}

\begin{itemize}
  \item Labels are case-sensitive and must be unique in the document.
  \item Prefixes such as \texttt{sec:}, \texttt{fig:}, \texttt{tab:},
    \texttt{eq:}, \texttt{thm:} keep them readable.
  \item After adding or moving a label, compile again (twice in total)
    so that the references are updated.
\end{itemize}

% ===============================================================
\section{Footnotes}
% ===============================================================

A footnote is written with \verb|\footnote{...}|, which numbers it
and places the text at the bottom of the page.%
\footnote{The footnote is typeset at the bottom of the page where it
appears, not at the end of the document.}
Inside a table or a \texttt{minipage}, use \verb|\footnotemark| in the
text and \verb|\footnotetext| after the table.

% ===============================================================
\section{Presentation conventions}
% ===============================================================

\subsection{Front matter and back matter}

A presentable document has, in order: a title block, an abstract, a
table of contents, lists of figures and tables, the body, and then
references and appendices. The sectioning commands fill the body:
\verb|\appendix| changes the later sections to letters.

\subsection{Spacing and emphasis}

\begin{itemize}
  \item Put one blank line between paragraphs; do not use \verb|\\| for
    that purpose.
  \item Use \verb|\emph| for the first use of a term and
    \verb|\textbf| sparingly.
  \item Keep display maths and figures close to the text that discusses
    them.
\end{itemize}

\appendix
\section{A short appendix}

Appendix sections are labelled A, B, \ldots. Here the theorem counter
continues to use the section number from the body, as
Theorem~\ref{thm:pyth} shows.

\end{document}
